\documentclass[11pt,a4paper]{article}

% ---------- packages ----------
\usepackage[utf8]{inputenc}
\usepackage[T1]{fontenc}
\usepackage{lmodern}
\usepackage{amsmath,amssymb}
\usepackage[margin=2.5cm]{geometry}
\usepackage{hyperref}
\usepackage{booktabs}
\usepackage{enumitem}
\usepackage{microtype}
\usepackage{parskip}
\usepackage{titlesec}
\usepackage{xcolor}

\hypersetup{
    colorlinks=true,
    linkcolor=blue!60!black,
    urlcolor=blue!60!black,
    citecolor=blue!60!black
}

\titleformat{\section}{\Large\bfseries\color{blue!60!black}}{\thesection}{0.7em}{}
\titleformat{\subsection}{\large\bfseries\color{blue!50!black}}{\thesubsection}{0.6em}{}

\newcommand{\R}{\mathbb{R}}
\newcommand{\vect}[1]{\mathbf{#1}}
\newcommand{\mat}[1]{\mathbf{#1}}
\newcommand{\E}{\mathbb{E}}
\newcommand{\Var}{\mathrm{Var}}
\newcommand{\Cov}{\mathrm{Cov}}

\title{\textbf{Multidimensional Statistical Methods (KMA/MSM)}\\[4pt]
\large A Plain-Language Guide for Computer Scientists}
\author{Compiled from the course materials of B.~Šedivá (ZČU)}
\date{}

\begin{document}

\maketitle
\tableofcontents
\newpage

% ======================================================================
\section*{Preface: How to Use This Guide}
\addcontentsline{toc}{section}{Preface}

This document is not a replacement for the lectures or the official
slides.  It is a \emph{translation} of the core ideas into the language
of someone who thinks in algorithms, data structures, and code.
The goal is to answer three questions for every method:

\begin{enumerate}[label=\textbf{\arabic*.}]
    \item \textbf{What is it?} (a plain description, no formulas first)
    \item \textbf{Why does it work?} (the geometric or probabilistic intuition)
    \item \textbf{Why do simpler alternatives fail?} (the trap the method avoids)
\end{enumerate}

Mathematics is included only where it sharpens the intuition.
You are not expected to reproduce proofs; you are expected to
understand \emph{why} a tool is the right one for a given job.

% ======================================================================
\section{The Big Picture: Why Multivariate Statistics?}

In univariate statistics you have \emph{one} number per observation
(height, salary, temperature).  In multivariate statistics you have a
\emph{vector} of numbers per observation:
\[
\vect{x}_i = (x_{i1}, x_{i2}, \dots, x_{ip}) \in \R^p .
\]
The dimensions are not independent.  A tall person tends to weigh more;
a high-income neighbourhood tends to have larger houses.  If you
analyse each variable separately, you throw away the information
hidden in the \emph{correlations}.  Worse, you create two problems:

\begin{itemize}
    \item \textbf{Multiple testing.}  If you run $p$ separate tests at
          significance level $0.05$, the chance of at least one false
          positive grows rapidly with $p$.
    \item \textbf{Geometry changes.}  In high dimensions, the Euclidean
          distance between points behaves counter-intuitively (the
          ``curse of dimensionality'').  A point can be far from the
          centre in one coordinate but still be a perfectly normal
          observation because the cloud is stretched in that direction.
\end{itemize}

Multivariate methods solve these problems by treating the whole vector
as one object.  The covariance matrix is the key that encodes the
shape of the cloud.

% ======================================================================
\section{Random Vectors and Multivariate Distributions}

\subsection{Joint, marginal, and conditional distributions}
A \textbf{random vector} $\vect{X}=(X_1,\dots,X_p)$ has a
\emph{joint distribution}.  The \textbf{marginal} distribution of
$X_1$ is what you get when you ignore all other coordinates.
The \textbf{conditional} distribution of $X_1$ given $X_2=x_2$ is what
you get when you \emph{fix} $X_2$ and look at the remaining
uncertainty.

\textbf{Why this matters for a computer scientist:}
A joint distribution is like a multidimensional array of probabilities.
Marginalisation is summing over an axis.  Conditioning is slicing the
array at a fixed index and renormalising.  Nothing mysterious—but the
notation in mathematical statistics often hides this simple picture.

\subsection{Covariance matrix}
The \textbf{covariance matrix} $\Sigma$ has entries
\[
\Sigma_{jk} = \Cov(X_j,X_k).
\]
The diagonal contains the variances; the off-diagonal contains the
correlations scaled by the standard deviations.
$\Sigma$ is symmetric and positive semi-definite.  It describes the
\emph{shape} of the probability cloud: the directions in which the
cloud is elongated or compressed.

\textbf{Intuition.}  Think of $\Sigma$ as the ``stretch matrix'' of an
ellipse.  If $\Sigma$ were the identity matrix, the cloud would be a
sphere.  If $\Sigma$ has a large eigenvalue in some direction, the
cloud is stretched along that direction.  This is exactly the
information PCA will exploit later.

% ======================================================================
\section{The Multivariate Normal Distribution}

\subsection{Why it is everywhere}
The multivariate normal (MVN) distribution is the natural
generalisation of the bell curve.  It is defined by a mean vector
$\boldsymbol{\mu}$ and a covariance matrix $\Sigma$.  Its density is
\[
f(\vect{x}) = \frac{1}{(2\pi)^{p/2}|\Sigma|^{1/2}}
\exp\!\left(-\frac12(\vect{x}-\boldsymbol{\mu})^\top
\Sigma^{-1}(\vect{x}-\boldsymbol{\mu})\right).
\]
The quadratic form inside the exponent,
\[
d_M^2(\vect{x}) = (\vect{x}-\boldsymbol{\mu})^\top
\Sigma^{-1}(\vect{x}-\boldsymbol{\mu}),
\]
is the \textbf{Mahalanobis distance}.  It measures how far a point is
from the centre \emph{in the metric of the data cloud}.

\textbf{Why it works.}  The level sets of the MVN density are
ellipsoids.  The Mahalanobis distance ``undoes'' the stretching of the
cloud, so that a point that lies along a stretched direction is not
penalised for being far away.  In contrast, Euclidean distance treats
all directions equally and will flag perfectly normal observations as
outliers simply because the data cloud is elongated.

\textbf{Why other distributions are harder.}  For non-normal data, the
covariance matrix does not fully describe the shape.  The MVN is
special because \emph{marginals and conditionals are again normal} and
linear transformations remain normal.  This closure property makes the
math tractable and is the reason most classical multivariate methods
assume (approximate) normality.

% ======================================================================
\section{Descriptive Statistics for Multivariate Data}

The sample mean is a vector:
\[
\bar{\vect{x}} = \frac1n \sum_{i=1}^n \vect{x}_i .
\]
The sample covariance matrix is
\[
\mat{S} = \frac{1}{n-1}\sum_{i=1}^n
(\vect{x}_i-\bar{\vect{x}})(\vect{x}_i-\bar{\vect{x}})^\top .
\]
The \textbf{correlation matrix} $\mat{R}$ is obtained by standardising
each variable (subtract mean, divide by standard deviation) before
computing the covariance.  It removes the units and makes the
correlations comparable.

\textbf{Practical note.}  Always standardise your data before PCA or
clustering unless the variables are already on the same scale.
Otherwise a variable measured in millimetres will dominate one
measured in metres, even if the latter is more informative.

% ======================================================================
\section{Outlier Detection}

\subsection{Euclidean distance is naive}
The Euclidean distance from the mean,
\[
d_E(\vect{x}) = \|\vect{x}-\bar{\vect{x}}\|,
\]
ignores the correlation structure.  A point that lies along the long
axis of the ellipse will have a large Euclidean distance even though
it is a typical observation for that direction.

\subsection{Mahalanobis distance}
The Mahalanobis distance
\[
d_M(\vect{x}) = \sqrt{(\vect{x}-\bar{\vect{x}})^\top
\mat{S}^{-1}(\vect{x}-\bar{\vect{x}})}
\]
accounts for the shape of the cloud.  It is the multivariate analogue
of a $z$-score.  Under the MVN assumption, $d_M^2$ follows a
chi-square distribution with $p$ degrees of freedom, which gives a
principled threshold for flagging outliers.

\textbf{Why it works.}  $\mat{S}^{-1}$ down-weights directions in which
the data naturally varies a lot and up-weights directions in which the
data is tightly clustered.  A point is an outlier only if it is far
away \emph{relative to the natural variability} of the cloud.

\textbf{Pitfall.}  If the data contains many outliers, $\mat{S}$ itself
is distorted.  Robust versions (e.g.\ MCD—minimum covariance
determinant) estimate the covariance from the clean majority before
computing distances.

% ======================================================================
\section{Hypothesis Testing in Multiple Dimensions}

\subsection{Why not run $p$ separate $t$-tests?}
If you test each coordinate separately, you face the \textbf{multiple
comparisons problem}.  With $p=10$ independent tests at $\alpha=0.05$,
the probability of at least one false rejection is
$1-0.95^{10}\approx 40\%$, not $5\%$.
Moreover, separate tests ignore the correlations between coordinates,
so they can miss a difference that is only visible when coordinates
are combined.

\subsection{Hotelling's $T^2$}
Hotelling's $T^2$ is the multivariate generalisation of the
$t$-statistic.  For a sample of vectors $\vect{x}_1,\dots,\vect{x}_n$
from an MVN distribution, it tests $H_0:\boldsymbol{\mu}=\boldsymbol{\mu}_0$.
The statistic
\[
T^2 = n(\bar{\vect{x}}-\boldsymbol{\mu}_0)^\top
\mat{S}^{-1}(\bar{\vect{x}}-\boldsymbol{\mu}_0)
\]
is a scalar that combines all coordinates in a way that respects
their covariance.  Under $H_0$ it is related to an $F$-distribution.

\textbf{Intuition.}  $T^2$ asks: \emph{How far is the sample mean from
the hypothesised mean, measured in units of the natural variability
of the mean vector?}  If the difference lies along a direction in
which the sample mean is very stable, even a small difference is
significant.  If the difference lies along a direction with huge
variability, a large difference may be noise.

% ======================================================================
\section{Comparing Multiple Groups: ANOVA and MANOVA}

\subsection{The problem with multiple $t$-tests}
If you have $k$ groups and you compare all pairs with $t$-tests, you
inflate the Type~I error rate.  With 5 groups there are 10 pairwise
tests; the chance of at least one false positive is much larger than
$0.05$.

\subsection{One-way ANOVA}
ANOVA compares the \emph{between-group variance} to the
\emph{within-group variance}:
\[
F = \frac{\text{variance of group means}}{\text{average variance inside groups}}.
\]
If the group means are genuinely different, the numerator is large
relative to the denominator.  ANOVA is an \emph{omnibus} test: it
tells you that \emph{some} difference exists, but not which groups
differ.  Post-hoc tests (Tukey, Bonferroni, etc.) then localise the
difference while controlling the family-wise error rate.

\subsection{MANOVA}
When each observation is a vector rather than a scalar, MANOVA
(multivariate ANOVA) compares group mean \emph{vectors}.  It uses
matrices of between-group and within-group sums of squares and
cross-products instead of scalars.  The intuition is the same:
\emph{Is the separation between group centres large compared with the
spread inside the groups?}

\textbf{Why this matters for a CS student.}  Think of ANOVA as a
statistical version of ``is the between-cluster distance large
compared with the within-cluster scatter?''.  It is the same idea
that underlies many clustering validity indices.

% ======================================================================
\section{Categorical Data and Contingency Tables}

Not all data are continuous.  A \textbf{contingency table} counts how
many observations fall into each combination of categories.
For two categorical variables with $r$ and $c$ categories, the table
has $r\times c$ cells.

The \textbf{chi-square test of independence} compares the observed
counts $O_{ij}$ with the counts $E_{ij}$ expected if the two variables
were independent:
\[
\chi^2 = \sum_{i,j}\frac{(O_{ij}-E_{ij})^2}{E_{ij}}.
\]
A large $\chi^2$ means the observed table is far from what independence
would predict.

\textbf{Intuition.}  If two variables are independent, the count in a
cell should be roughly (row total $\times$ column total) / grand total.
The chi-square statistic aggregates the squared relative deviations.
It is the categorical analogue of ``correlation''.

\textbf{Multivariate extension.}  For more than two categorical
variables, log-linear models are the standard tool.  They model the
log of the expected cell counts as a sum of main effects and
interactions, exactly like ANOVA but for counts.

% ======================================================================
\section{Correlation and Its Pitfalls}

\subsection{Pearson correlation}
The Pearson correlation $\rho_{XY}$ measures the strength of a
\emph{linear} relationship between two variables.  In a multivariate
setting, the correlation matrix is the standardised version of the
covariance matrix.

\subsection{Partial correlation}
A high correlation between $X$ and $Y$ may be \textbf{spurious}: both
may be driven by a third variable $Z$.  The \textbf{partial
correlation} $\rho_{XY\cdot Z}$ measures the correlation between $X$
and $Y$ after removing the linear effect of $Z$.
If the partial correlation is much smaller than the raw correlation,
the raw correlation was largely an artefact of $Z$.

\textbf{Why this matters.}  In causal inference, confusing correlation
with causation is the classic trap.  Partial correlation is the
simplest way to ``control for'' a confounder.

\subsection{Copulas (briefly)}
A \textbf{copula} separates the marginal distributions from the
dependence structure.  It says: ``You can have any marginal
distributions you like, and a separate copula that glues them together
with a chosen dependence pattern.''  This is useful when the
normal assumption fails but you still want to model dependence.

% ======================================================================
\section{Regression with Multiple Predictors}

\subsection{Multiple linear regression}
You model a response $Y$ as a linear combination of predictors
$\vect{X}=(X_1,\dots,X_p)$ plus noise:
\[
Y = \beta_0 + \beta_1 X_1 + \cdots + \beta_p X_p + \varepsilon .
\]
The coefficients are estimated by least squares: minimise the sum of
squared residuals.  In matrix form, $\hat{\boldsymbol{\beta}} =
(\mat{X}^\top\mat{X})^{-1}\mat{X}^\top\vect{y}$.

\textbf{Why it works.}  Least squares is the maximum-likelihood
estimate under the assumption that the noise is normal with constant
variance.  It is also the best linear unbiased estimator (BLUE) under
weaker assumptions (Gauss–Markov theorem).

\subsection{Multicollinearity}
When predictors are highly correlated, $\mat{X}^\top\mat{X}$ becomes
nearly singular.  The coefficient estimates become unstable: small
changes in the data cause huge changes in $\hat{\boldsymbol{\beta}}$.
This is \textbf{multicollinearity}.

\textbf{Why it matters.}  A model with unstable coefficients may
predict well on the training data but generalise poorly.  It also
makes interpretation impossible: you cannot say ``increasing $X_1$ by
one unit increases $Y$ by $\beta_1$'' if $X_1$ and $X_2$ always move
together.

\subsection{Regularisation: Ridge and Lasso}
To combat multicollinearity and overfitting, add a penalty to the
least-squares objective:
\[
\text{Ridge: } \|\vect{y}-\mat{X}\boldsymbol{\beta}\|^2 +
\lambda\|\boldsymbol{\beta}\|_2^2,
\qquad
\text{Lasso: } \|\vect{y}-\mat{X}\boldsymbol{\beta}\|^2 +
\lambda\|\boldsymbol{\beta}\|_1 .
\]
Ridge shrinks coefficients towards zero; Lasso can set some
coefficients exactly to zero, performing automatic feature selection.

\textbf{Intuition.}  The penalty is a budget on the size of the
coefficients.  It prevents the model from relying too heavily on any
single predictor, which is especially important when predictors are
correlated.

% ======================================================================
\section{Dimension Reduction: Principal Component Analysis (PCA)}

\subsection{The problem}
You have $p$ variables, many of which are correlated.  You suspect that
the data lie close to a lower-dimensional subspace, but you cannot
see it because of the noise and the coordinate system.

\subsection{What PCA does}
PCA finds a new coordinate system in which the axes are ordered by
the amount of variance they explain.  The first principal component
(PC1) is the direction of maximum variance.  PC2 is the direction of
maximum variance \emph{orthogonal} to PC1, and so on.

Mathematically, the principal components are the eigenvectors of the
covariance matrix $\Sigma$ (or correlation matrix $\mat{R}$), ordered
by their eigenvalues.  The eigenvalue $\lambda_j$ is the variance
explained by PC $j$.

\subsection{Why it works}
The covariance matrix encodes the shape of the data cloud.  Its
eigenvectors are the axes of the ellipsoid that best fits the cloud;
the eigenvalues are the lengths of those axes.  By projecting onto
the top $k$ eigenvectors, you keep the directions where the data
varies most and discard the directions where it is almost constant.

\textbf{Why not just pick random dimensions?}  Because the original
coordinates are arbitrary.  A variable that looks important in the
original space may actually be a combination of several underlying
factors.  PCA finds the \emph{natural} axes of the data, independent
of the original measurement units (provided you standardise first).

\subsection{Pitfalls}
\begin{itemize}
    \item \textbf{Scale matters.}  Always standardise unless all
          variables are in the same units.
    \item \textbf{PCA is unsupervised.}  It does not know about class
          labels.  The directions of maximum variance may not be the
          directions that separate classes.
    \item \textbf{Interpretation is hard.}  A principal component is a
          linear combination of all original variables; it rarely has
          a simple name.
\end{itemize}

% ======================================================================
\section{Discriminant Analysis: LDA and QDA}

\subsection{The goal}
You have labelled data: each observation belongs to one of $K$
classes.  You want to find a rule that assigns new observations to
the correct class.

\subsection{Linear Discriminant Analysis (LDA)}
LDA assumes that each class has a multivariate normal distribution
with the \emph{same} covariance matrix $\Sigma$, but different means
$\boldsymbol{\mu}_k$.  It finds linear combinations of the variables
that maximise the ratio
\[
\frac{\text{between-class variance}}{\text{within-class variance}}.
\]
In other words, it looks for directions in which the class centres
are far apart relative to the spread inside each class.

\textbf{Why LDA is not PCA.}  PCA finds directions of maximum overall
variance, ignoring labels.  LDA finds directions that are \emph{useful
for classification}.  A direction with huge variance may be useless
for separating classes if all classes vary in the same way along that
direction.  Conversely, a direction with small variance may be
perfectly separating the classes.

\subsection{Quadratic Discriminant Analysis (QDA)}
QDA relaxes the assumption of equal covariance matrices: each class
gets its own $\Sigma_k$.  The decision boundaries become quadratic
instead of linear.  QDA is more flexible but needs more data to
estimate the extra parameters.  With small samples, LDA is often more
stable even if its assumptions are slightly violated.

\subsection{PCA--LDA}
A common pipeline is to first reduce dimensions with PCA (to denoise
and compress), then apply LDA on the principal components.  This works
well when the data is high-dimensional and the classes are linearly
separable after compression, but it can also hurt if the PCA step
discards directions that were important for classification.

% ======================================================================
\section{Clustering}

\subsection{The goal}
Clustering is unsupervised: you have no labels, but you want to group
observations so that points in the same group are similar and points
in different groups are dissimilar.

\subsection{K-means}
K-means partitions the data into $K$ clusters by alternating between
two steps:
\begin{enumerate}
    \item Assign each point to the nearest cluster centre.
    \item Recompute each centre as the mean of the points assigned to it.
\end{enumerate}
It minimises the within-cluster sum of squares.

\textbf{Why it works.}  It is a greedy algorithm for a non-convex
optimisation problem.  It converges to a local optimum, which may
depend on the initialisation.  K-means++ is a smarter initialisation
that often finds better local optima.

\subsection{Hierarchical clustering}
Hierarchical clustering builds a tree (dendrogram) of nested clusters.
It does not require you to choose $K$ in advance; you can cut the
dendrogram at a height that gives a sensible number of clusters.

\textbf{Distance matters.}  Euclidean distance in the original space
can be misleading if variables are correlated or on different scales.
Standardise first, or use Mahalanobis distance.

% ======================================================================
\section{Generating Multivariate Data}

Why would a computer scientist care about generating multivariate
data?  For simulation, testing algorithms, and bootstrapping.

The standard trick for generating MVN data is the \textbf{Cholesky
decomposition}.  If $\Sigma = \mat{L}\mat{L}^\top$ with $\mat{L}$
lower triangular, then
\[
\vect{X} = \boldsymbol{\mu} + \mat{L}\vect{Z},
\qquad \vect{Z}\sim\mathcal{N}(\vect{0},\mat{I}),
\]
has covariance matrix $\Sigma$.

\textbf{Intuition.}  Start with independent standard normal variables
(a spherical cloud).  Multiply by $\mat{L}$ to stretch and rotate the
cloud into the desired shape.  The Cholesky factor is the ``square
root'' of the covariance matrix, analogous to taking the square root
of the variance in the univariate case.

% ======================================================================
\section{Putting It All Together: A Decision Guide}

\begin{center}
\begin{tabular}{@{}p{0.28\textwidth}p{0.32\textwidth}p{0.32\textwidth}@{}}
\toprule
\textbf{Goal} & \textbf{Method} & \textbf{Key assumption / caveat} \\
\midrule
Summarise a multivariate cloud & Mean vector, covariance matrix &
Sensitive to outliers; use robust versions if needed. \\
Detect outliers & Mahalanobis distance &
Needs a good estimate of $\Sigma$; robust MCD helps. \\
Test a mean vector & Hotelling's $T^2$ &
Multivariate normality; large $n$ helps. \\
Compare $k$ group means & ANOVA / MANOVA &
Independence, normality, homogeneity of variance. \\
Analyse categorical counts & Chi-square / log-linear models &
Expected counts should be large enough. \\
Measure dependence & Correlation, partial correlation &
Linear relationships only; watch for confounders. \\
Predict a continuous response & Multiple regression, Ridge/Lasso &
Check multicollinearity; regularise if $p$ is large. \\
Reduce dimensions & PCA &
Standardise first; unsupervised, may discard class-relevant directions. \\
Classify labelled data & LDA / QDA &
MVN assumption; LDA needs equal covariances. \\
Find natural groups & K-means, hierarchical clustering &
Choose $K$ or cut height; distance metric matters. \\
\bottomrule
\end{tabular}
\end{center}

% ======================================================================
\section*{Appendix: A Minimal Mathematical Toolkit}
\addcontentsline{toc}{section}{Appendix: A Minimal Mathematical Toolkit}

\begin{itemize}
    \item \textbf{Vector:} an ordered list of numbers, $\vect{x}\in\R^p$.
    \item \textbf{Matrix:} a rectangular array; $\mat{A}^\top$ is the transpose.
    \item \textbf{Eigenvector/eigenvalue:} $\mat{A}\vect{v}=\lambda\vect{v}$.
          Eigenvectors are directions that are only stretched (not rotated)
          by the matrix; eigenvalues are the stretch factors.
    \item \textbf{Positive semi-definite:} a symmetric matrix $\mat{A}$
          such that $\vect{x}^\top\mat{A}\vect{x}\ge0$ for all $\vect{x}$.
          Covariance matrices are always positive semi-definite.
    \item \textbf{Cholesky decomposition:} $\Sigma=\mat{L}\mat{L}^\top$
          with $\mat{L}$ lower triangular.  The multivariate analogue of
          a square root.
\end{itemize}

\vspace{1cm}
\noindent\rule{\textwidth}{0.4pt}
\footnotesize
This guide is an independent explanation written for learning purposes.
It is not an official document of ZČU and does not replace the
lectures, slides, or the teacher’s own materials.
\normalsize

\end{document}
